(1) 梁弯曲应变能 \(U_b=(EI/2)\int_0^L(w'')^2dx\)，轴向恒压缩死载因几何缩短而做功，对应势能贡献 \(U_P=-(P/2)\int_0^L(w')^2dx\)。总二阶势能 \(\Pi_2=U_b+U_P\)。变分并两次分部积分得到内部 Euler--Lagrange 方程 \(EIw''''+Pw''=0\)。铰支位移固定 \(w(0)=w(L)=0\)，自然弯矩条件 \(EIw''(0)=EIw''(L)=0\)，排除了固支边界的错误基频。

(2) 取满足铰支条件的 \(w_n=A_n\sin(n\pi x/L)\)（整数 \(n\ge1\)）。代入平衡方程，得到 \([EI(n\pi/L)^4-P(n\pi/L)^2]A_n=0\)，非零振幅要求 \(P_n=EI(n\pi/L)^2\)。这些特征载荷按 \(n^2\) 单调增加，最先出现的第一失稳为 \(P_{\rm cr}=\pi^2EI/L^2\)，对应正弦单半波形状。高阶模态虽也是形式本征解，但不能在达到第一临界之前由零态自发失稳。

(3) 稳定直梁需要二阶势能对所有非零允许扰动严格为正。由 Fourier 正弦级数或 Poincaré 不等式，对于满足端位移条件的 \(w\)，有 \(\int_0^L(w'')^2dx\ge(\pi/L)^2\int_0^L(w')^2dx\)。因此 \(\Pi_2\ge[EI(\pi/L)^2-P]\int_0^L(w')^2dx/2\)。若 \(P<P_{\rm cr}\) 所有非零扰动均使势能增加；若 \(P>P_{\rm cr}\)，选择第一正弦模可使 \(\Pi_2<0\)。于是 Rayleigh 商 \(EI\int(w'')^2/\int(w')^2\) 的极小值恰为 \(P_{\rm cr}\)。

(4) 令 \(w=A\sin(\pi x/L)\)，直接积分给 \(\Pi_2(A)=A^2L[EI\pi^4/L^4-P\pi^2/L^2]/4=\pi^2A^2[P_{\rm cr}-P]/(4L)\)。当 \(P=P_{\rm cr}\)，二阶项恰为零并不能断言任意振幅皆为平衡；当 \(P>P_{\rm cr}\)，二阶理论仅说明直梁失稳，不能预测屈曲后的幅值或后屈曲分支。稳定有限挠度依赖几何非线性、轴向伸缩与载荷控制方式，高阶能量必须重新建立。
